\begin{definition}[One-break Herbrand functions]\label{mod:ramification-herbrand}
Let $t\in\mathbf N$ be the unique lower break from
equation~\ref{eq:one-break-new}, and let $\ell>0$ be the degree.  Define
total real functions
\[
 \varphi_{t,\ell}(u)
 =\min\{u,t\}+\frac{\max\{u-t,0\}}{\ell},
 \qquad
 \psi_{t,\ell}(v)
 =\min\{v,t\}+\ell\max\{v-t,0\}.
\]
On the manuscript's domain $[-1,\infty)$, the first is the specialization
of the integral in equation~\ref{eq:Herbrand-def}: the one-break density is
$1$ through $t$ and $1/\ell$ strictly after $t$.  The same formulas give a
harmless total extension below $-1$.

The exact closed forms, with both clauses including the break, are
\[
 \varphi_{t,\ell}(u)=
 \begin{cases}
  u,&u\le t,\\
  t+(u-t)/\ell,&t\le u,
 \end{cases}
 \qquad
 \psi_{t,\ell}(v)=
 \begin{cases}
  v,&v\le t,\\
  t+\ell(v-t),&t\le v.
 \end{cases}
\]
Thus the two descriptions agree at the endpoint and
\[
 \varphi_{t,\ell}(-1)=\psi_{t,\ell}(-1)=-1,
 \qquad
 \varphi_{t,\ell}(t)=\psi_{t,\ell}(t)=t.
\]
Both functions are continuous and monotone.  When $\ell>0$, they are
strictly increasing and mutually inverse in the two separately stated
directions
\[
 \psi_{t,\ell}(\varphi_{t,\ell}(u))=u,
 \qquad
 \varphi_{t,\ell}(\psi_{t,\ell}(v))=v.
\]
Moreover,
\(\varphi_{t,\ell}(u)\le u\le\psi_{t,\ell}(u)\).
For an actual \texttt{PrimeCyclicExtension} and a supplied
\texttt{IsLowerBreak} witness, the declaration \texttt{herbrand} applies
these inverse identities with the canonical degree
$\ell=[K:F]$; primality supplies its positivity.  No formula is stored as
a structure field.

For the integral unit-filtration depths in Theorems~\ref{thm:graded-norm}
and~\ref{thm:norm-filtration}, define
\[
 \Psi_{t,\ell}(r)=\min\{r,t\}+\ell(r-t)
 =\begin{cases}
   r,&r\le t,\\
   t+\ell(r-t),&t\le r,
  \end{cases}
 \qquad(r\in\mathbf N),
\]
where the subtraction on the first line is natural-number subtraction.
Then
\[
 (\Psi_{t,\ell}(r):\mathbf R)=\psi_{t,\ell}(r),
 \qquad
 \varphi_{t,\ell}(\Psi_{t,\ell}(r))=r.
\]
Consequently the real floor and ceiling of $\psi_{t,\ell}(r)$ both recover
the same integral depth.  The converter is strictly increasing, monotone,
tends to infinity, and satisfies $r\le\Psi_{t,\ell}(r)$.  At and above the
break,
\[
 \Psi_{t,\ell}(t)=t,
 \qquad
 \Psi_{t,\ell}(t+1)=t+\ell,
 \qquad
 \Psi_{t,\ell}(r+k)=\Psi_{t,\ell}(r)+\ell k.
\]
In particular,
\(\Psi_{t,\ell}(r+1)=\Psi_{t,\ell}(r)+\ell\).
The second source exponent in the norm theorem is exactly the natural
depth $\Psi_{t,\ell}(r)+1$.  When $\ell$ is prime,
\[
 \Psi_{t,\ell}(r)+1<\Psi_{t,\ell}(r+1)
 \qquad(t\le r),
\]
so that off-by-one exponent is not confused with the next Herbrand level.
The remaining integer conversions used by the later stationary-depth
arguments are
\[
 \Psi_{t,\ell}\!\left(\left\lfloor\frac r2\right\rfloor\right)
 \le
 \left\lfloor\frac{\Psi_{t,\ell}(r)}2\right\rfloor
\]
and, for $t\le s$,
\[
 \Psi_{t,\ell}(2s+e)
 =2\Psi_{t,\ell}(s)+(\ell-1)t+\ell e.
\]
The latter is proved for every $e\in\mathbf N$, although the manuscript
uses only $e\in\{0,1\}$.  Finally, if $t+2\le d$, then the exact
subtraction-free form of the high stationary-depth calculation is
\[
 \Psi_{t,\ell}(2d+\varepsilon-1)+1
 =2\bigl(\Psi_{t,\ell}(d-1)+1\bigr)
   +(\ell-1)(t+1)+\ell\varepsilon.
\]
Again Lean proves this for every natural $\varepsilon$, hence in particular
for the two parity values used in the manuscript.

\LeanFile{LanglandsFirstMainLemma/Ramification/Herbrand.lean}
\lean{LanglandsFirstMainLemma.herbrandPhi}
\lean{LanglandsFirstMainLemma.herbrandPsi}
\lean{LanglandsFirstMainLemma.herbrandPsiNat}
\lean{LanglandsFirstMainLemma.herbrandPhi_of_le_break}
\lean{LanglandsFirstMainLemma.herbrandPhi_of_break_le}
\lean{LanglandsFirstMainLemma.herbrandPsi_of_le_break}
\lean{LanglandsFirstMainLemma.herbrandPsi_of_break_le}
\lean{LanglandsFirstMainLemma.herbrandPhi_break}
\lean{LanglandsFirstMainLemma.herbrandPsi_break}
\lean{LanglandsFirstMainLemma.herbrandPhi_neg_one}
\lean{LanglandsFirstMainLemma.herbrandPsi_neg_one}
\lean{LanglandsFirstMainLemma.herbrandPhi_monotone}
\lean{LanglandsFirstMainLemma.herbrandPsi_monotone}
\lean{LanglandsFirstMainLemma.continuous_herbrandPhi}
\lean{LanglandsFirstMainLemma.continuous_herbrandPsi}
\lean{LanglandsFirstMainLemma.herbrandPsi_herbrandPhi}
\lean{LanglandsFirstMainLemma.herbrandPhi_herbrandPsi}
\lean{LanglandsFirstMainLemma.herbrandPsi_leftInverse_herbrandPhi}
\lean{LanglandsFirstMainLemma.herbrandPsi_rightInverse_herbrandPhi}
\lean{LanglandsFirstMainLemma.herbrandPhi_strictMono}
\lean{LanglandsFirstMainLemma.herbrandPsi_strictMono}
\lean{LanglandsFirstMainLemma.herbrandPhi_le_self}
\lean{LanglandsFirstMainLemma.self_le_herbrandPsi}
\lean{LanglandsFirstMainLemma.herbrandPsiNat_of_le_break}
\lean{LanglandsFirstMainLemma.herbrandPsiNat_of_break_le}
\lean{LanglandsFirstMainLemma.herbrandPsiNat_break}
\lean{LanglandsFirstMainLemma.herbrandPsiNat_cast}
\lean{LanglandsFirstMainLemma.herbrandPhi_herbrandPsiNat_cast}
\lean{LanglandsFirstMainLemma.ceil_herbrandPsi_natCast}
\lean{LanglandsFirstMainLemma.floor_herbrandPsi_natCast}
\lean{LanglandsFirstMainLemma.herbrandPsi_natCast_add_one}
\lean{LanglandsFirstMainLemma.herbrandPsiNat_add}
\lean{LanglandsFirstMainLemma.herbrandPsiNat_succ}
\lean{LanglandsFirstMainLemma.herbrandPsiNat_break_succ}
\lean{LanglandsFirstMainLemma.herbrandPsiNat_strictMono}
\lean{LanglandsFirstMainLemma.herbrandPsiNat_monotone}
\lean{LanglandsFirstMainLemma.tendsto_herbrandPsiNat_atTop}
\lean{LanglandsFirstMainLemma.self_le_herbrandPsiNat}
\lean{LanglandsFirstMainLemma.herbrandPsiNat_add_one_lt_succ}
\lean{LanglandsFirstMainLemma.herbrandPsiNat_two_mul_add}
\lean{LanglandsFirstMainLemma.herbrandPsiNat_half_le_half}
\lean{LanglandsFirstMainLemma.herbrandPsiNat_stationary_depth}
\lean{LanglandsFirstMainLemma.herbrand}
\leanok
\SourceText{Equations~\ref{eq:Herbrand-def}, \ref{eq:one-break-new}, and
\ref{eq:Herbrand-prime-new}; Theorems~\ref{thm:graded-norm} and
\ref{thm:norm-filtration}; and Lemma~\ref{lem:stationary-depth-norm}.}
\uses{mod:ramification-primecyclicextension}
\FormalizationContract
\end{definition}
